\documentclass[10pt,a4paper]{amsart}
\usepackage{amsmath,amssymb,amsthm,graphicx}
% For commutative diagrams we use Xy-pic
\usepackage[all]{xy}

% We use 2cell for 2-commutative diagrams.
\xyoption{2cell}
% Theorem environments.
%
\theoremstyle{plain}
\newtheorem{theorem}[subsection]{Theorem}
\newtheorem{proposition}[subsection]{Proposition}
\newtheorem{lemma}[subsection]{Lemma}

\theoremstyle{definition}
\newtheorem{definition}[subsection]{Definition}
\newtheorem{example}[subsection]{Example}
\newtheorem{exercise}[subsection]{Exercise}
\newtheorem{situation}[subsection]{Situation}

\theoremstyle{remark}
\newtheorem{remark}[subsection]{Remark}
\newtheorem{remarks}[subsection]{Remarks}

\numberwithin{equation}{subsection}

% Macros
%
\def\lim{\mathop{\mathrm{lim}}\nolimits}
\def\colim{\mathop{\mathrm{colim}}\nolimits}
\def\Spec{\mathop{\mathrm{Spec}}}
\def\Hom{\mathop{\mathrm{Hom}}\nolimits}
\def\Ext{\mathop{\mathrm{Ext}}\nolimits}
\def\SheafHom{\mathop{\mathcal{H}\!\mathit{om}}\nolimits}
\def\SheafExt{\mathop{\mathcal{E}\!\mathit{xt}}\nolimits}
\def\Sch{\mathit{Sch}}
\def\Mor{\mathop{\mathrm{Mor}}\nolimits}
\def\Ob{\mathop{\mathrm{Ob}}\nolimits}
\def\Sh{\mathop{\mathit{Sh}}\nolimits}
\def\NL{\mathop{N\!L}\nolimits}
\def\CH{\mathop{\mathrm{CH}}\nolimits}
\def\proetale{{pro\text{-}\acute{e}tale}}
\def\etale{{\acute{e}tale}}
\def\QCoh{\mathit{QCoh}}
\def\Ker{\mathop{\mathrm{Ker}}}
\def\Im{\mathop{\mathrm{Im}}}
\def\Coker{\mathop{\mathrm{Coker}}}
\def\Coim{\mathop{\mathrm{Coim}}}

% Boxtimes
%
\DeclareMathSymbol{\boxtimes}{\mathbin}{AMSa}{"02}

%
% Macros for moduli stacks/spaces
%
\def\QCohstack{\mathcal{QC}\!\mathit{oh}}
\def\Cohstack{\mathcal{C}\!\mathit{oh}}
\def\Spacesstack{\mathcal{S}\!\mathit{paces}}
\def\Quotfunctor{\mathrm{Quot}}
\def\Hilbfunctor{\mathrm{Hilb}}
\def\Curvesstack{\mathcal{C}\!\mathit{urves}}
\def\Polarizedstack{\mathcal{P}\!\mathit{olarized}}
\def\Complexesstack{\mathcal{C}\!\mathit{omplexes}}
% \Pic is the operator that assigns to X its picard group, usage \Pic(X)
% \Picardstack_{X/B} denotes the Picard stack of X over B
% \Picardfunctor_{X/B} denotes the Picard functor of X over B
\def\Pic{\mathop{\mathrm{Pic}}\nolimits}
\def\Picardstack{\mathcal{P}\!\mathit{ic}}
\def\Picardfunctor{\mathrm{Pic}}
\def\Deformationcategory{\mathcal{D}\!\mathit{ef}}
\newenvironment{reference}{}{}
\newenvironment{slogan}{}{}
\newenvironment{history}{}{}
\usepackage{cleveref}
\crefname{originallemma}{original lemma}{original lemmas}
\Crefname{originallemma}{Original lemma}{Original lemmas}
\makeatletter
\newlabel{lemma-closure-is-model}{{0C2S}{source}{}{sourcepointer-1}{}}
\newlabel{lemma-closure-is-model@cref}{{[originallemma][0][]0C2S}{[1][0][]source}{}{}{}}
\makeatother
\crefname{originallemma}{original lemma}{original lemmas}
\Crefname{originallemma}{Original lemma}{Original lemmas}
\makeatletter
\newlabel{lemma-regular}{{0C2U}{source}{}{sourcepointer-2}{}}
\newlabel{lemma-regular@cref}{{[originallemma][0][]0C2U}{[1][0][]source}{}{}{}}
\makeatother
\crefname{originallemma}{original lemma}{original lemmas}
\Crefname{originallemma}{Original lemma}{Original lemmas}
\makeatletter
\newlabel{lemma-pre-exists-minimal-model}{{0CD9}{source}{}{sourcepointer-3}{}}
\newlabel{lemma-pre-exists-minimal-model@cref}{{[originallemma][0][]0CD9}{[1][0][]source}{}{}{}}
\makeatother
\title{Existence of a minimal model for a smooth projective curve over the fraction field of a discrete valuation ring}
\author{The Stacks Project Authors}
\date{stacks-project:models.tex:local-snapshot-20261012; GNU FDL 1.2 or later}
\begin{document}
\maketitle
\noindent\textit{Editorial note.} The counted unit is the Stacks Project Proposition proposition-exists-minimal-model at source lines 4518-4522 together with its proof at source lines 4524-4534, a straight adjacent pairing in the chapter file. The chapter is the upstream file models.tex, which is the file name recorded in the source record; the pool directory carries the longer name stacks-models and is not used as the chapter name anywhere in this card.\par\medskip
\noindent\textit{Editorial note.} This chapter's candidate index records the statement body and the proof body without their closing markers, so the retained spans used here are the full environments, exactly as in the earlier cards of this chapter.\par\medskip
\noindent\textit{Editorial note.} The retained proof names three results of the same chapter, for the scheme theoretic image being a model, for the existence of a resolution of singularities, and for the construction of a minimal model from a regular proper one. All three are declared as non-load-bearing navigation pointers with their stable Stacks tags looked up in the shared tag index. The counted argument is the two-step reduction itself.\par\medskip
\noindent\textit{Editorial note.} The statement mentions no situation and no labelled object, so no context span is retained, and no citation occurs in the retained spans, so no bibliography entry is supplied.\par\medskip
\noindent\textit{Editorial note.} The retained preamble is the chapter preamble's Xy-pic setup at lines 14-18, the theorem declarations and mathematical macros at lines 154-217, and the three no-op wrapper environments.\par\medskip
\noindent\textit{Allowed standard prerequisites (editorial).}\begin{itemize}
\item Models, resolutions of singularities in dimension one, and the pre-existence lemma. A model of a smooth projective curve over the fraction field of a discrete valuation ring is a flat proper scheme over the ring with that curve as generic fibre; the chapter proves separately that a certain construction always produces a model that admits a morphism from every other regular proper model, and it records the resolution of singularities available for the one-dimensional schemes that occur here. The author's own statements of those results are the allowed inputs.. The counted proposition is the existence statement on which the definition of a minimal model and the uniqueness results of the surrounding sections rest.
\end{itemize}\medskip
\ifdefined\hypertarget\hypertarget{sourcepointer-1}{}\fi
\noindent\textit{Source reference.} Original lemma 0C2S: focused-seeds/AGP/sources/stacks-snapshot-20261011/tags.txt, source lines 15167--15167. Navigation pointer invoked by the retained proof to the effect that the scheme theoretic image of the curve is a projective model; the cited statement is an allowed input and is not reproduced inside the span. The number is the stable Stacks tag of the target, looked up in the shared tag index.\par\medskip
\ifdefined\hypertarget\hypertarget{sourcepointer-2}{}\fi
\noindent\textit{Source reference.} Original lemma 0C2U: focused-seeds/AGP/sources/stacks-snapshot-20261011/tags.txt, source lines 15169--15169. Navigation pointer invoked by the retained proof for the resolution of singularities of the projective model; the cited existence statement is an allowed input and is not reproduced inside the span. The number is the stable Stacks tag of the target, looked up in the shared tag index.\par\medskip
\ifdefined\hypertarget\hypertarget{sourcepointer-3}{}\fi
\noindent\textit{Source reference.} Original lemma 0CD9: focused-seeds/AGP/sources/stacks-snapshot-20261011/tags.txt, source lines 15534--15534. Navigation pointer invoked by the retained proof for the construction that turns a regular proper model into a minimal one; the cited statement is an allowed input and is not reproduced inside the span. The number is the stable Stacks tag of the target, looked up in the shared tag index.\par\medskip
\begin{proposition}
\label{proposition-exists-minimal-model}
Let $C$ be a smooth projective curve over $K$ with
$H^0(C, \mathcal{O}_C) = K$. A minimal model exists.
\end{proposition}

\begin{proof}
Choose a closed immersion $C \to \mathbf{P}^n_K$. Let
$X$ be the scheme theoretic image of $C \to \mathbf{P}^n_R$.
Then $X \to \Spec(R)$ is a projective model of $C$ by
Lemma \ref{lemma-closure-is-model}.
By Lemma \ref{lemma-regular} there exists a resolution
of singularities $X' \to X$ and $X'$ is a model for $C$.
Then $X' \to \Spec(R)$ is proper as a composition of proper morphisms.
Then we may apply Lemma \ref{lemma-pre-exists-minimal-model}
to obtain a minimal model.
\end{proof}

\end{document}
